% !TeX root = ../../main.tex

Die oberen Abschnitte haben bisher gezeigt, wie Datenbanken im Allgemeinen Fehler finden und behandeln.
Des Weiteren zeigen sie, wie der Forschungsstand zu den Fehlerklassen in Bezug auf Retry aussieht.
Und wie schon erwähnt, bringt Retry nicht nur Vorteile, sondern auch Nachteile für die Performance des Systems.
Deshalb beschreibt dieser Abschnitt, wie sich Retry auf ein System auswirkt.

Die Arbeit von Aderaldo und Mendonça beschreibt das Retry Verhalten auf ein System passend.
In Ihrer Arbeit testen sie wie sich verschiedene Retry Strategien und Resilienzbibliotheken auf das System auswirken.
Als flexible Parameter haben sie \emph{Initial backoff delay}, \emph{Backoff delay multiplier} und
\emph{Max. number of retries} definiert.
Die beiden Resilienzbibliotheken die getestet wurden, waren Polly(C\#) und Resilience4j(Java).
Sie testen anhand der Werte \emph{Relative Execution Time}, \emph{Relative Response Time} und
\emph{Relative Performance}.
Es wurden vier Forschungsfragen gestellt.
\begin{itemize}
	\item Wie ist die execution time unter den verschiedenen Retry Konfigurationen, Bibliotheken, Workloads und
	Failure Rates?
	\item Wie ist die response time unter den verschiedenen Retry Konfigurationen, Bibliotheken, Workloads und
	Failure Rates?
	\item Welche Konfiguration bietet das beste Ergebnis?
	\item Wie wirken sich die Parameter aus?
\end{itemize}
Um die ersten beiden Fragen zu beantworten, haben Sie jeweils mit nur einem variablen Parameter gearbeitet.
Die anderen Beiden blieben fix.
Für die dritte Frage haben Sie die fünf besten Retry Konfigurationen für die beiden Bibliotheken unter den
unterschiedlichen Workloads und Failure Rates dargestellt.
Da Sie für die vierte Frage \emph{SHAP}\footnote{SHAP ist ein machine learning explainer. Mit dem werden in der Regel
	erklärt, wie ein machine learning model auf seine Ergebnisse kommt.} verwendet haben, wurden die Ergebnisse mit
einem SHAP Value für die einzelnen Parameter bewertet und nach den Kombinationen aus Workloads und Failure Rates
sortiert.

Aus diesen Ergebnissen lassen sich die Forschungsfragen beantworten.
Die execution time steigt exponentiell an, wenn die einzelnen Parameter und die failure rate höher werden.
Während der backoff delay multiplier hier einen stärkeren Einfluss hat als die restlichen Parameter, hat Workload
hier keinen signifikanten Einfluss.
Die Ergebnisse für die zweite Frage zeigen, die response time geht runter.
Dies wird wieder am stärksten von dem backoff delay multiplier beeinflusst.
Hier sorgen höhere Parameter für eine geringere response time, besonders unter hohen failure rates.
Es lässt sich nur schwer festlegen, welche Konfiguration die Beste ist, da die Parameter die execution und response
time ausbalancieren.
Durch höhere Parameter steigt die execution time, aber die response time sinkt.
Gleichzeitig geht die Optimierung der einen Metrik meist zu Lasten der anderen, sobald die failure rates steigen.
Vor allem da die Einstellungen keinem festen Muster folgen, vor allem unter hoher Workload und failure rates.
Dadurch ist es hart die beste Strategie zu bestimmen, aber es gibt sie.
An dem Punkt an dem sowohl die execution time, als auch die response time im Verhältnis zueinander am niedrigsten
sind, befindet sich die beste Strategie.

Als Fazit ziehen Aderaldo und Mendonça, dass eine feine Balance der Parameter notwendig ist, damit das System
sauber funktioniert.
Deshalb ist es unter anderem wichtig die Betriebsumgebung zu berücksichtigen, damit man die Parameter passend
abstimmen kann~\cite{Aderaldo2023}.